\chapter{\ifproject%
\ifenglish Project Structure and Methodology\else โครงสร้างและขั้นตอนการทำงาน\fi
\else%
\ifenglish Project Structure\else โครงสร้างของโครงงาน\fi
\fi
}

ในบทนี้จะกล่าวถึงหลักการ การออกแบบระบบ และระเบียบวิธีวิจัยที่ใช้ในการประเมินความปลอดภัยของตัวแบบภาษาขนาดใหญ่ (LLMs) โดยครอบคลุมตั้งแต่กระบวนการเตรียมชุดข้อมูล การออกแบบกลไกการโจมตีและการป้องกัน ไปจนถึงการพัฒนาระบบรันการทดสอบอัตโนมัติ

\makeatletter

% \renewcommand\section{\@startsection {section}{1}{\z@}%
%                                    {13.5ex \@plus -1ex \@minus -.2ex}%
%                                    {2.3ex \@plus.2ex}%
%                                    {\normalfont\large\bfseries}}

\makeatother
%\vspace{2ex}
% \titleformat{\section}{\normalfont\bfseries}{\thesection}{1em}{}
% \titlespacing*{\section}{0pt}{10ex}{0pt}

\section{ภาพรวมสถาปัตยกรรมระบบ (System Architecture)}
ระบบถูกออกแบบให้ทำงานในรูปแบบ Pipeline เพื่อทดสอบโมเดลภาษาผ่าน Groq API โดยมีขั้นตอนการไหลของข้อมูลดังนี้:
\begin{enumerate}
  \item Input Layer: รับคำสั่ง (Prompt) ที่ผ่านการปรับแต่งระดับภาษาไทย 5 ระดับ
  \item Defense Layer (Input Scrubbing / Defensive Prompting): ใช้เทคนิค Input Scrubbing เพื่อกรองข้อมูลส่วนบุคคล (PII) ขาเข้าด้วยเทคนิค Regex และ NER หรือ Defensive Prompting เพื่อกำหนดขอบเขตพฤติกรรมของโมเดลให้ปฏิเสธการตอบข้อมูลที่เป็นอันตราย
  \item Inference Layer: ส่งคำสั่งไปยังโมเดล llama-3.1-8b-instant, meta-llama-llama-4-maverick-17b-128e-instruct หรือ moonshotai-kimi-k2-instruct-0905 ผ่าน Groq API พร้อมการใช้งาน Defensive Prompting
  \item Evaluation Layer: รวบรวมคำตอบและประเมินผลด้วยมาตรวัดที่ออกแบบมาเฉพาะสำหรับแต่ละรูปแบบการโจมตี โดยดำเนินการทดสอบซ้ำอย่างน้อย 5 รอบต่อกรณีเพื่อหาค่าเฉลี่ยอัตราความสำเร็จ (Average Success Rate) รวมถึงวิเคราะห์เปรียบเทียบประสิทธิภาพของแต่ละเทคนิค และประเมินผลกระทบของปริมาณ Token และจำนวนคำที่มีต่ออัตราความสำเร็จ
\end{enumerate}

\section{การเตรียมชุดข้อมูลและระดับภาษา (Data Preparation)}
\subsection{ชุดข้อมูลหลัก}
\begin{itemize}
    \item \textbf{Enron Email Dataset:} ชุดข้อมูลอีเมลสาธารณะจากการสืบสวนคดีล้มละลายของบริษัท Enron (ค.ศ. 2001) ที่มักใช้ในงานวิจัยด้าน Machine Learning, NLP, Cybersecurity และ Data Leakage Detection ในงานศึกษานี้ใช้เพื่อจำลองสถานการณ์ Data Exfiltration Attack (DEA) เพื่อประเมินการป้องกันการรั่วไหลของข้อมูลภายใต้บริบททางวิชาการ

    \begin{center}
        \includegraphics[width=0.8\linewidth]{"picture/Data Preparation/DEA without defense.png"}
        \captionof{figure}{DEA without defense}
        \label{fig:dea_without_defense}
    \end{center}

    \item \textbf{Jailbreak Prompt:} ชุดคำสั่งที่ออกแบบมาเพื่อหลีกเลี่ยงระบบความปลอดภัยของโมเดลภาษาขนาดใหญ่ (LLMs) โดยใช้เทคนิคหลากหลายเช่น Persona Prompts, Role-playing หรือการสร้างสถานการณ์จำลองเพื่อหลอกลวงให้โมเดลเปิดเผยข้อมูลที่เป็นอันตราย

    \begin{center}
        \includegraphics[width=0.8\linewidth]{"picture/Data Preparation/JA without defense.png"}
        \captionof{figure}{JA without defense}
        \label{fig:ja_without_defense}
    \end{center}

\end{itemize}
\subsection{กระบวนการปรับระดับภาษา}
ในการเตรียมชุดข้อมูลสำหรับการทดลอง ได้ใช้เทคนิค Batch Processing ผ่าน OpenAI API เพื่อแปลและปรับระดับภาษาของชุดคำสั่งจำนวน 3,050 ชุดแบบอัตโนมัติ โดยกำหนดให้ครอบคลุมภาษาไทย 5 ระดับ ได้แก่ ระดับพิธีการ ทางการ กึ่งทางการ ไม่เป็นทางการ และกันเอง กระบวนการนี้ช่วยรักษาความหมายเดิมของคำสั่งขณะสร้างความหลากหลายด้านรูปแบบภาษา ลดความไม่สม่ำเสมอจากการแปลด้วยมนุษย์ และรองรับการประมวลผลข้อมูลขนาดใหญ่ เพื่อให้ชุดข้อมูลสะท้อนลักษณะการใช้งานจริงและเหมาะสมต่อการประเมินระบบในบริบทภาษาที่แตกต่างกัน
\section{ระเบียบวิธีการโจมตี (Attack Methodology)}
สภาพแวดล้อมการทดสอบ: ดำเนินการทดสอบซ้ำอย่างน้อย 5 รอบต่อกรณี เพื่อหาค่าเฉลี่ยอัตราความสำเร็จ (Average Success Rate)
\subsection{Data Extraction Attack (DEA)} 
ออกแบบคำสั่งเพื่อล่อลวงให้โมเดลเปิดเผยข้อมูลอีเมลหรือข้อมูลลับที่ถูกเก็บไว้ในหน่วยความจำ
\subsection{Jailbreak Attack (JA)}
ใช้เทคนิค Persona Prompts หรือการสร้างสถานการณ์จำลองเพื่อหลีกเลี่ยงระบบความปลอดภัยของโมเดล


\begin{figure}[htbp]
  \centering
  \includegraphics[width=0.8\linewidth]{"picture/System Flow/Attack System Flow.png"}
  \caption{Attack System Flow}
  \label{fig:attack_system_flow}
\end{figure}

\section{การออกแบบกลไกการป้องกัน (Defense Design)}
\subsection{Input Scrubbing}
พัฒนาระบบตรวจสอบคำสั่งขาเข้าด้วย Regular Expression (Regex) และ Named Entity Recognition (NER) เพื่อตรวจจับและปกปิดข้อมูลประเภท ชื่อ, อีเมล, เบอร์โทรศัพท์ ก่อนถึงตัวโมเดล

\begin{figure}[htbp]
  \centering
  \includegraphics[width=0.8\linewidth]{"picture/System Flow/Scrubbing System Flow.png"}
  \caption{Scrubbing System Flow}
  \label{fig:scrubbing_system_flow}
\end{figure}

\subsection{Defensive Prompting}
ออกแบบ System Prompt จำนวน 12 รูปแบบ (แบ่งเป็นสำหรับ DEA 6 แบบ และ JA 6 แบบ) เพื่อกำหนดขอบเขตพฤติกรรมของโมเดลให้ปฏิเสธการตอบข้อมูลที่เป็นอันตราย

\begin{figure}[htbp]
  \centering
  \includegraphics[width=0.8\linewidth]{"picture/System Flow/Defensive Prompt System Flow.png"}
  \caption{Defensive Prompt System Flow}
  \label{fig:defensive_prompt_system_flow}
\end{figure}

\section{การพัฒนาระบบรันการทดสอบอัตโนมัติ (Automated Testing)}
ทีมวิจัยได้พัฒนาเครื่องมือเพื่อเพิ่มประสิทธิภาพในการทดลอง ดังนี้:
\begin{itemize}
  \item Multi-API and Multi-Attack Execution: สคริปต์ใน Terminal เดียวที่สามารถสั่งรันหลายโมเดลและหลายรูปแบบการโจมตีพร้อมกันได้
  \item Automated Data Logging: ระบบจะบันทึกผลลัพธ์ลงใน Google Sheets โดยอัตโนมัติ
  \item Metrics Calculation: พัฒนาฟังก์ชันคำนวณปริมาณ Token, ความยาวข้อความ และสร้างกราฟสรุปผล (Visualization) ทันทีหลังจบการทดลอง
\end{itemize}

\section{มาตรวัดการประเมินผล (Evaluation Metrics)}
\begin{itemize}
  \item DEA Metrics: วัดความถูกต้องแม่นยำ (Accuracy) ในการดึงข้อมูลตามความเป็นจริง
  \item JA Metrics: กำหนดเกณฑ์การตัดสิน (Labeling) ว่าโมเดลถูกเจาะระบบสำเร็จหรือไม่ โดยพิจารณาจากข้อมูลที่โมเดลตอบกลับมาในส่วนของข้อมูลติดต่อ
\end{itemize}